% ATTEMPT_STATUS: unresolved
% TOP_PROBLEM: {"problemId": "problem.non-perturbative-construction-and-non-triviality-of-4d-euclidean-phi4-quantum-field-theory", "canonicalTitle": "Non-Perturbative Construction and Non-Triviality of 4D Euclidean Phi4 Quantum Field Theory", "releaseRank": 51, "primaryDomain": "mathematical_physics", "primaryDomainLabel": "Mathematical physics", "displayStatus": "reviewed_hold", "releaseStatus": "open"}
% !TeX program = lualatex
% Complete problem section copied verbatim from research_batch_51_54.tex.
\documentclass[11pt]{article}
\usepackage[margin=25mm]{geometry}
\usepackage{fontspec}
\setmainfont{Latin Modern Roman}
\usepackage{amsmath,amssymb,mathtools}
\usepackage{unicode-math}
\setmathfont{Latin Modern Math}
\usepackage{enumitem,needspace}
\usepackage{xurl,hyperref}
\hypersetup{colorlinks=true,urlcolor=blue}
\setcounter{secnumdepth}{1}
\setlength{\parindent}{0pt}
\setlength{\parskip}{5pt}
\setlength{\emergencystretch}{3em}
\setlist[itemize]{leftmargin=3em,itemsep=5pt,topsep=4pt,parsep=0pt}
\allowdisplaybreaks
\newcommand{\N}{\mathbb{N}}
\newcommand{\Z}{\mathbb{Z}}
\newcommand{\Q}{\mathbb{Q}}
\newcommand{\R}{\mathbb{R}}
\newcommand{\C}{\mathbb{C}}
\newcommand{\F}{\mathbb{F}}
\newcommand{\A}{\mathbb{A}}
\newcommand{\PP}{\mathbb P}
\newcommand{\E}{\mathbb{E}}
\newcommand{\eps}{\varepsilon}
\newcommand{\dd}{\,\mathrm d}
\newcommand{\sref}[2]{\hyperlink{source:#1:#2}{[S#2]}}
\newcommand{\eref}[2]{\hyperlink{extra:#1:#2}{[E#2]}}
% Turn the next switch off for a shorter reading copy. The quotations preserve
% every exactTarget field, including qualifications omitted from a short summary.
\newif\ifshowcatalogue
\showcataloguetrue
\newcommand{\cataloguescope}[1]{\ifshowcatalogue
 \paragraph{Exact scope in the supplied catalogue.}
 \begin{quote}\small #1\end{quote}\fi}
\begin{document}
\setcounter{section}{50}
% ================================================================
% problemId: problem.non-perturbative-construction-and-non-triviality-of-4d-euclidean-phi4-quantum-field-theory
% source releaseRank: 51; source releaseStatus: open
% exactTarget SHA-256: f31f95cf191e6998f9d82a2d67885af59277b6579c9e6ee2e8b6d57c561713eb
\clearpage
\section{\texorpdfstring{Non-Perturbative Construction and Non-Triviality of 4D Euclidean Phi4 Quantum Field Theory}{Non-Perturbative Construction and Non-Triviality of 4D Euclidean Phi4 Quantum Field Theory}}\label{51}
\subsection{Definitions and mathematical statement}
A Euclidean scalar field is a probability law on generalized functions on $\R^4$, or equivalently a compatible collection of Schwinger correlation distributions. The formal quartic action is
\[
 S(\phi)=\int_{\R^4}\left(\tfrac12|\nabla\phi|^2+\tfrac12m^2\phi^2+\lambda\phi^4\right)dx.
\]
This expression by itself does not define the continuum measure. A construction must specify regularization, renormalization, and a limit satisfying the Osterwalder--Schrader conditions, including Euclidean invariance, reflection positivity, symmetry, and regularity. Non-Gaussian means that the correlations are not all determined by the two-point function through Wick's rule. \emph{Source-fit warning:} the cited Aizenman--Duminil-Copin theorem proves Gaussian limits for its four-dimensional ferromagnetic lattice Ising/$\phi^4$ class. The requested non-Gaussian construction cannot use that class of limits; the catalogue does not specify an alternative class. The record is retained, not certified as a well-posed currently open construction problem.
\par\smallskip\textit{Formulation and concept references:} \sref{51}{1}, \eref{51}{1}.
\cataloguescope{Construct a non-trivial, non-Gaussian Euclidean scalar quantum field theory with quartic interaction (phi\textasciicircum{}4)\_4 in four spacetime dimensions satisfying the Osterwalder-Schrader axioms.}
\subsection{Short English statement}
Can the proposed four-dimensional quartic scalar theory be constructed as a genuinely non-Gaussian continuum field, outside the class already ruled out by the cited triviality theorem?
% BEGIN RESEARCH ADDITION ATTEMPT-51
\subsection{Research attempt}

\paragraph{Current frontier.}
The first obstruction is a theorem, not merely perturbative evidence.
Aizenman--Duminil-Copin's Theorem~1.2 applies to the scaling-limit
construction in their Section~1.2, with its Griffiths--Simon single-site
measures, ferromagnetic nearest-neighbour interaction, normalization and
critical/near-critical hypotheses. Within that construction the limiting
field is Gaussian. This does not quantify over every possible meaning of a
continuum quartic theory. \eref{51}{1}
Panis extends quantitative Gaussianity to specified long-range models with
effective dimension at least four; his assumptions (A1)--(A6) and the
four-dimensional conclusion in Corollary~1.14 matter. Merely replacing
nearest-neighbour interactions by a longer-range ferromagnetic kernel is
therefore not an unrestricted escape route. \eref{51}{2}

Two catalogue references require a different epistemic status. Manolessou's
Theorem~3.2 claims reflection positivity for a previously constructed
hierarchy; that earlier construction and the full continuum package have
not been independently verified here. It is not used as an established
solution. \sref{51}{3}, \eref{51}{5}
Borji's August 2026 analysis concerns a closed second-order mean-field
hierarchy; its concluding discussion explicitly leaves the ultraviolet
remainder outside the result. It supplies neither a general construction
nor a universal no-go theorem. \sref{51}{4}, \eref{51}{4}

\paragraph{Attack route.}
Three routes were compared. Staying in the standard lattice class would
require contradicting the cited Gaussianity theorem, so that route is
excluded under its hypotheses. A positive covariance with stronger
ultraviolet decay might make a quartic measure easier to construct, but it
must retain reflection positivity. A cutoff-dependent perturbation of a
known Gaussian-limit measure might instead retain positivity while changing
the limit; that requires a perturbation that remains visible on continuum
observables. The calculations below investigate the latter two routes and
then test a tempting non-Gaussian construction based on mixing free fields.
They are deductions made for this notebook, not claims of literature novelty.

\paragraph{Attempt: a spectral obstruction to ultraviolet softening.}
\textbf{Lemma 51.1 (positive spectral mass cannot disappear at high
momentum).} Suppose a scalar covariance multiplier has the unsubtracted
representation
\begin{equation}\label{ra:51:stieltjes}
 c(s)=\int_{[0,\infty)}\frac{\rho(\dd t)}{s+t},\qquad s>0,
 \qquad 0\ne\rho\ge0,
\end{equation}
and the integral is finite for each $s>0$. Then
\[
 \lim_{s\to\infty}s\,c(s)=\rho([0,\infty))\in(0,\infty].
\]
Indeed, $s/(s+t)$ increases to one at every finite $t$, so the monotone
convergence theorem applies, including when the total mass is infinite.
In particular, $c(s)=o(s^{-1})$ is impossible under
\eqref{ra:51:stieltjes}. The representation is an explicit hypothesis:
this argument does not assume that every distributional covariance admits
an unsubtracted representation of this form.

Consider the concrete higher-derivative proposal
\[
 c_\varepsilon(s)=\frac{1}{s+m^2+\varepsilon s^2},\qquad
 m>0,\quad 0<4\varepsilon m^2<1.
\]
Put
\[
 a=\frac{1-\sqrt{1-4\varepsilon m^2}}{2\varepsilon},\qquad
 b=\frac{1+\sqrt{1-4\varepsilon m^2}}{2\varepsilon}.
\]
Then $0<a<b$ and the exact partial-fraction identity is
\begin{equation}\label{ra:51:ghost}
 c_\varepsilon(s)=\frac{1}{\varepsilon(b-a)}
 \left(\frac{1}{s+a}-\frac{1}{s+b}\right).
\end{equation}
The negative residue is consistent with the general rational-propagator
criterion of Arici et al., Theorem~3.7. Here it is useful to exhibit the
negative reflection form directly rather than merely invoke the criterion.
\eref{51}{3}

Reflect the first Euclidean coordinate, and Fourier transform in the
remaining three. For fixed spatial momentum $k$, set
$\omega_a=(|k|^2+a)^{1/2}$, $\omega_b=(|k|^2+b)^{1/2}$, and
$\mathcal Lf(\omega)=\int_0^\infty e^{-\omega t}f(t)\dd t$.
With the usual inverse Fourier normalization, the reflected quadratic form
for a temporal test function is
\begin{equation}\label{ra:51:reflection}
 Q_k(f)=\frac{1}{\varepsilon(b-a)}
 \left(\frac{|\mathcal Lf(\omega_a)|^2}{2\omega_a}
       -\frac{|\mathcal Lf(\omega_b)|^2}{2\omega_b}\right).
\end{equation}
This follows by inverse transforming $(p_0^2+\omega^2)^{-1}$, whose temporal
kernel is $e^{-\omega|t|}/(2\omega)$; reflection makes the argument $t+s$.

Choose a nonnegative nonzero $h\in C_c^\infty(0,\infty)$ and $0<t_1<t_2$.
For a fixed $k_0$, take
\[
 f(t)=h(t-t_1)-e^{\omega_a(k_0)(t_2-t_1)}h(t-t_2).
\]
Its Laplace transform vanishes at $\omega_a(k_0)$ and is nonzero at
$\omega_b(k_0)$, since these frequencies are different. Thus
$Q_{k_0}(f)<0$. By continuity the inequality persists in a neighbourhood
of $k_0$. Multiplying by a nonzero smooth spatial Fourier test function
supported in that neighbourhood and integrating gives a negative
reflection form for a Schwartz test function supported at positive time.
Therefore this particular free regulator fails reflection positivity at
every indicated finite $\varepsilon$, despite its positive Fourier
multiplier on real momenta. No numerical approximation is involved.

This rejects a proposed positive-metric regulator, not the target theory.
As $\varepsilon\downarrow0$, the heavy pole moves to infinity; reflection
positivity might conceivably be restored in a limit. An interacting
construction using such regulators would still have to prove tightness,
renormalized convergence, limiting reflection positivity and a nonzero
connected correlation. The finite-cutoff calculation establishes none of
those limiting assertions.

\paragraph{Attempt: how large must a perturbation be to escape Gaussianity?}
\textbf{Lemma 51.2 (change-of-measure bound).} Let $\mu_a$ be cutoff
probability measures and let $\nu_a=r_a\mu_a$, where $r_a\ge0$ and
$\int r_a\dd\mu_a=1$. Fix $1<p<\infty$, $q=p/(p-1)$, and write
$\delta_a=\|r_a-1\|_{L^p(\mu_a)}$. Let $Y_{1,a},\ldots,Y_{4,a}$ be real
observables, centred under both measures, satisfying
\[
 \max_i\|Y_{i,a}\|_{L^{4q}(\mu_a)}\le M
\]
with $M$ independent of $a$. Their connected fourth moments satisfy
\begin{equation}\label{ra:51:cumulantbound}
 |\kappa_{4,\nu_a}(Y_{1,a},\ldots,Y_{4,a})
       -\kappa_{4,\mu_a}(Y_{1,a},\ldots,Y_{4,a})|
 \le M^4(7\delta_a+3\delta_a^2).
\end{equation}

To prove this, H\"older's inequality gives
\[
 |\E_{\nu_a}Y_1Y_2Y_3Y_4-\E_{\mu_a}Y_1Y_2Y_3Y_4|
 \le\delta_a M^4,
 \qquad
 |\E_{\nu_a}Y_iY_j-\E_{\mu_a}Y_iY_j|
 \le\delta_a M^2.
\]
The second estimate also uses monotonicity of $L^r$ norms on a probability
space. Every reference covariance has absolute value at most $M^2$.
Writing a perturbed covariance as $C_{ij}+\Delta_{ij}$, the difference of
each product of two covariances is at most
$M^4(2\delta_a+\delta_a^2)$. Subtract the three pairing products from the
fourth moment. This proves \eqref{ra:51:cumulantbound}. Centring is not
silently inherited after tilting: for example, it is guaranteed when both
measures are invariant under field sign reversal and all $Y_i$ are odd.
Without such a condition additional mean terms must be included.

There is also a stronger qualitative observation that needs no moment
bound. For any finite vector $Y_a$ of smeared observables and $t\in\R^j$,
\begin{equation}\label{ra:51:characteristic}
 \left|\E_{\nu_a}e^{i t\cdot Y_a}-\E_{\mu_a}e^{i t\cdot Y_a}\right|
 \le\|r_a-1\|_{L^1(\mu_a)}\le\delta_a.
\end{equation}
Consequently, if $\delta_a\to0$, every existing Gaussian finite-dimensional
limit under $\mu_a$ remains the same under $\nu_a$. This statement uses
bounded characteristic functions, not an unjustified inference from weak
convergence to convergence of fourth moments. Applied to the lattice
limits of \eref{51}{1}, it says that asymptotically negligible density
tilts cannot provide the requested non-Gaussian limit.

For a quantitative witness, suppose the reference fourth cumulant tends
to zero and the absolute limiting perturbed cumulant is at least $\eta>0$.
Under the uniform moment hypothesis, \eqref{ra:51:cumulantbound} forces
\[
 \liminf_{a\to0}\delta_a
 \ge \frac{\sqrt{49+12\eta/M^4}-7}{6}>0.
\]
This is a necessary condition, not a sufficient construction criterion.
A small bare coefficient in $r_a=Z_a^{-1}e^{-V_a}$ does not imply
$\delta_a\to0$: the growing volume, ultraviolet fluctuations and the
normalizing partition function all intervene. Conversely, the lemma does
not rule out a perturbation large in microscopic total variation but small
on a particular finite set of smeared observables.

\paragraph{Stress test: a non-Gaussian field that does not solve the problem.}
Let $\Psi$ be a centred massive free Euclidean field and let $V$ be an
independent positive, bounded, nonconstant random variable with $\E V=1$.
The field $\Phi=\sqrt V\,\Psi$ is non-Gaussian. Conditional Wick expansion
gives, for $C_{ij}=\E\Psi(f_i)\Psi(f_j)$,
\[
 \kappa_4(\Phi(f_1),\ldots,\Phi(f_4))
 =\operatorname{Var}(V)
   (C_{12}C_{34}+C_{13}C_{24}+C_{14}C_{23}).
\]
Euclidean invariance and reflection positivity survive the mixture:
the reflected expectation of $\overline{F(\theta\Phi)}F(\Phi)$ is an
average of nonnegative free-field reflected expectations. This elementary
construction therefore defeats the false claim that reflection positivity
alone forces Gaussianity.

It also exposes why those checks are insufficient. If $g_R$ is a large
translate of $g$, the massive free cross-covariance tends to zero, whereas
\[
 \operatorname{Cov}\bigl(\Phi(f)^2,\Phi(g_R)^2\bigr)
 \longrightarrow
 \operatorname{Var}(V)\,C(f,f)\,C(g,g)>0.
\]
Thus the mixture has a persistent global random scale and fails the
clustering requirement in the usual vacuum formulation. Independently of
which clustering convention is intended, no local quartic interaction has
been exhibited. The mixture is not a candidate answer to Section~\ref{51}.

The two substantive failed routes are now precise. Stronger ultraviolet
falloff conflicts with a positive unsubtracted spectral measure, and the
explicit higher-derivative covariance has a negative reflected direction.
Small measure tilts preserve a Gaussian scaling limit under the stated
hypotheses. Neither argument covers all local non-Gaussian continuum
constructions. In particular, absence of a fourth cumulant would not by
itself prove an arbitrary law Gaussian; the full characteristic-function
argument above is what establishes Gaussianity for negligible tilts.

\paragraph{Outcome.}
\textbf{Outcome: Exploratory approach with unresolved gap.}

The obtained results are the explicit negative-reflection test, the
quantitative fourth-cumulant stability bound, and the diagnosis of the
Gaussian-mixture false construction. They constrain proposed escapes from
the established lattice triviality theorem but do not produce a quartic
continuum field or rule out every such field. A successful continuation
must specify an alternative regularized model, retain an order-one
non-Gaussian effect on continuum observables, and prove all required
continuum axioms. The unspecified alternative model in the original entry
remains a genuine formulation limitation, not a licence to substitute an
arbitrary non-Gaussian random field.
% END RESEARCH ADDITION ATTEMPT-51
\subsection{Sources}
\begin{itemize}\raggedright
\item[\textnormal{[S1]}]\hypertarget{source:51:1}{} M. Aizenman, H. Duminil-Copin, Ann. of Math. 194(1):163-235, 2021.
\url{https://doi.org/10.4007/annals.2021.194.1.3}.
{\small\textit{Catalogue formulation source.}}
\item[\textnormal{[S2]}]\hypertarget{source:51:2}{} Marginal triviality of the scaling limits of critical 4D Ising and phi4 models.
\url{https://annals.math.princeton.edu/2021/194-1/p03}.
{\small\textit{Catalogue research/status reference; not a proof endorsement.}}
\item[\textnormal{[S3]}]\hypertarget{source:51:3}{} The non triviality of a Phi4 model, III, the Osterwalder-Schrader Positivity.
\url{https://arxiv.org/abs/2101.11980}.
{\small\textit{Catalogue research/status reference; not a proof endorsement.}}
\item[\textnormal{[S4]}]\hypertarget{source:51:4}{} Triviality in a Non-Perturbative Second-Order Mean-Field Theory for phi4.
\url{https://arxiv.org/abs/2608.05998}.
{\small\textit{Catalogue research/status reference; not a proof endorsement.}}
\item[\textnormal{[E1]}]\hypertarget{extra:51:1}{} Michael Aizenman and Hugo Duminil-Copin, Marginal triviality of the scaling limits of critical 4D Ising and phi4 models, Annals of Mathematics 194 (2021), 163--235; arXiv:1912.07973.
\url{https://arxiv.org/abs/1912.07973}.
{\small\textit{Added primary source: Source-fit warning and proved triviality for the paper's model class. Consulted 24 September 2026.}}
% BEGIN RESEARCH ADDITION REFERENCES-51
\item[\textnormal{[E2]}]\hypertarget{extra:51:2}{} Romain Panis, \emph{Triviality of the scaling limits of critical Ising and $\phi^4$ models with effective dimension at least four}, arXiv:2309.05797; Annals of Probability, DOI 10.1214/25-AOP1782. Assumptions (A1)--(A6), Theorems 1.12--1.13 and Corollary 1.14.
\url{https://arxiv.org/abs/2309.05797}.
{\small\textit{Added research source: Quantitative triviality for the specified long-range model class; not a no-go theorem for arbitrary continuum constructions. Consulted 24 September 2026.}}
\item[\textnormal{[E3]}]\hypertarget{extra:51:3}{} Francesca Arici, Daniel Becker, Chris Ripken, Frank Saueressig and Walter D.\ van Suijlekom, \emph{Reflection positivity in higher derivative scalar theories}, arXiv:1712.04308 (2017), Theorem 3.7; DOI 10.1063/1.5027231.
\url{https://arxiv.org/abs/1712.04308}.
{\small\textit{Added research source: Rational-covariance reflection-positivity criterion. The negative reflection form used here is also derived explicitly in the attempt. Consulted 24 September 2026.}}
\item[\textnormal{[E4]}]\hypertarget{extra:51:4}{} Majdouline Borji, \emph{Triviality in a Non-Perturbative Second-Order Mean-Field Theory for $\phi^4$}, arXiv:2608.05998 (August 2026), concluding discussion of the ultraviolet remainder.
\url{https://arxiv.org/abs/2608.05998}.
{\small\textit{Added research source: Mean-field result and its stated limitation; no assertion of full-theory triviality is imported. Consulted 24 September 2026.}}
\item[\textnormal{[E5]}]\hypertarget{extra:51:5}{} Marietta Manolessou, \emph{The ``non triviality'' of a $\Phi^4_4$ model, III, the ``Osterwalder--Schrader Positivity''}, arXiv:2101.11980 (2021), Theorem 3.2.
\url{https://arxiv.org/abs/2101.11980}.
{\small\textit{Added research source: Claimed positivity result. The preceding hierarchy construction and the full construction claim were not independently verified in this notebook; not a proof endorsement. Consulted 24 September 2026.}}
% END RESEARCH ADDITION REFERENCES-51
\end{itemize}

\end{document}
